% ============================================================================
%  ECE 431 -- Sheet 3 Solutions
%  Style: sheet-answering-style.md
% ============================================================================

\begin{question}[Question 1]If the lower energy $E_{1}$ of the two-energies model introduced in
lectures is the ground state then
\begin{enumerate}[label=(\alph*)]
  \item Re-write the rate equations given in lectures showing that they reduce to a single rate
        equation, and
  \item Solve this rate equation together with $N_{1}+N_{2}=N_{a}$, where $N_{a}$ is the
        atomic density, to find expressions of the population inversion density $N$ under both
        saturation and un-saturation conditions.
\end{enumerate}
\end{question}

\textbf{\large Solution --- Question 1}

\textbf{Part (a): Reduction to a single rate equation.}

When $E_1$ is the ground state, its lifetime is infinite ($\tau_1\to\infty$), so ground-state pump terms vanish ($R_1=0$). The full two-level rate equations become:

\begin{equation*}
  \frac{dN_2}{dt} = R + \sigma\Phi N_1 - \sigma\Phi N_2 - \frac{N_2}{\tau_2}, \qquad
  \frac{dN_1}{dt} = -R - \sigma\Phi N_1 + \sigma\Phi N_2 + \frac{N_2}{\tau_2} = -\frac{dN_2}{dt}.
\end{equation*}

The two equations satisfy conservation of mass $N_1+N_2=N_a$, so $dN_1/dt+dN_2/dt=0$, meaning \emph{only one equation is independent}. Substituting $N_1 = N_a - N_2$:

\begin{equation*}
  \boxed{\frac{dN_2}{dt} = R + \sigma\Phi N_a - 2\sigma\Phi N_2 - \frac{N_2}{\tau_2}.}
\end{equation*}

\textbf{Part (b): Population inversion density.}

At steady state ($dN_2/dt = 0$), with $N = N_2 - N_1 = 2N_2 - N_a$.

\emph{Un-saturation ($\Phi\to 0$):} the stimulated terms vanish, giving $N_2 = \tau_2 R$, so:

\begin{equation*}
  \boxed{N_o = 2\tau_2 R - N_a.}
\end{equation*}

\emph{Saturation ($\Phi\neq 0$):} solving $0 = R + \sigma\Phi N_a - (2\sigma\Phi + 1/\tau_2)N_2$:

\begin{equation*}
  N_2 = \frac{R + \sigma\Phi N_a}{\tfrac{1}{\tau_2}+2\sigma\Phi}
      = \frac{\tau_2 R + \tau_2\sigma\Phi N_a}{1+2\tau_2\sigma\Phi}.
\end{equation*}

Then $N = 2N_2 - N_a = \dfrac{2\tau_2 R + 2\tau_2\sigma\Phi N_a - N_a - 2\tau_2\sigma\Phi N_a}{1+2\tau_2\sigma\Phi}$:

\begin{equation*}
  \boxed{N = \frac{N_o}{1+\Phi/\Phi_s}, \qquad \Phi_s = \frac{1}{2\tau_2\sigma}.}
\end{equation*}

The inversion saturates with a characteristic saturation flux $\Phi_s$ inversely proportional to the stimulated-emission cross section and the upper-level lifetime.

% ============================================================================

\begin{question}[Question 2]A hypothetical gas laser uses Neon and an atomic transition at
$\lambda_{o}=0.5\,\mu$m that has a radiative decay time constant $\tau_{r}=10$\,ns, and a
transition cross section of $3\times10^{-16}$\,cm$^{-2}$. If the gas is at 300\,K then
calculate the FWHM susceptibility linewidth for the partial contribution of radiative and Doppler
broadening.
\end{question}

\textbf{\large Solution --- Question 2}

\textbf{Radiative (natural) linewidth.}

The natural linewidth is set by the radiative lifetime alone:

\begin{equation*}
  \Delta\nu_\text{rad} = \frac{1}{2\pi\tau_r} = \frac{1}{2\pi\times10^{-8}}
  \approx \boxed{15.62\;\text{MHz.}}
\end{equation*}

\textbf{Doppler linewidth.}

The given cross section is the Gaussian Doppler-broadened cross section at resonance:
$\sigma(\nu_o) = (\lambda^{2}/8\pi\tau_r)\times g_D(\nu_o)$, where $g_D(\nu_o)=1/(\sqrt{2\pi}\,\sigma_D)$. Solving for $\sigma_D$:

\begin{equation*}
  3\times10^{-16} = \frac{(0.5\times10^{-4})^{2}}{8\pi\times10^{-8}}\times\frac{1}{\sqrt{2\pi}\,\sigma_D}
  \quad\Longrightarrow\quad
  \sigma_D \approx 13.23\times10^{12}\;\text{Hz.}
\end{equation*}

The Doppler FWHM is therefore:

\begin{equation*}
  \boxed{\Delta\nu_\text{Dop} = \sqrt{8\ln 2}\,\sigma_D \approx 31.15\times10^{12}\;\text{Hz} = 31.15\;\text{THz.}}
\end{equation*}

Doppler broadening is three orders of magnitude wider than the natural linewidth, confirming that inhomogeneous broadening dominates completely in this gas laser at room temperature.

% ============================================================================

\begin{question}[Question 3]Use the gain coefficient formula derived in class to calculate the
$e$-folding distance $(1/\gamma)$ for intensity growth for the following laser:
$\lambda_{o}=0.6\,\mu$m, $N_{2}=6\times10^{9}$\,cm$^{-3}$, $N_{1}\approx 0$,
$g(\nu)\approx1/\Delta\nu$ where $\Delta\nu=0.5$\,GHz, $\tau_{r}=100$\,ns, and refractive
index $n\approx 1.0$.
\end{question}

\textbf{\large Solution --- Question 3}

The gain coefficient is $\gamma = \sigma N$ where $\sigma = (\lambda^{2}/8\pi\tau_r)\,g(\nu)$. With $g(\nu)\approx1/\Delta\nu$:

\begin{equation*}
  \gamma = \frac{\lambda^{2}}{8\pi\tau_r}\cdot\frac{1}{\Delta\nu}\cdot N_2
  = \frac{(0.6\times10^{-4})^{2}}{8\pi\times100\times10^{-9}}\times\frac{1}{0.5\times10^{9}}\times6\times10^{9}
  = 0.017\;\text{cm}^{-1}.
\end{equation*}

\begin{equation*}
  \boxed{\frac{1}{\gamma} = 58.17\;\text{cm.}}
\end{equation*}

The $e$-folding length of nearly 60\,cm shows how weakly the field amplifies in a single pass through a gas laser with only $6\times10^9$\,cm$^{-3}$ inverted atoms.

% ============================================================================

\begin{question}[Question 4]A commercially available ruby laser amplifier using a 15-cm-long rod has a
small-signal gain of 12. What is the small-signal gain of a 20-cm-long rod? Neglect gain
saturation effects.
\end{question}

\textbf{\large Solution --- Question 4}

For a laser amplifier the small-signal gain is $G_o = e^{\gamma_o L}$. Using the 15-cm rod to find $\gamma_o$:

\begin{equation*}
  12 = e^{\gamma_o\times15} \quad\Longrightarrow\quad \gamma_o = \frac{\ln 12}{15} = 0.165\;\text{cm}^{-1}.
\end{equation*}

For a 20-cm rod with the same gain coefficient:

\begin{equation*}
  \boxed{G_o\big|_{20\,\text{cm}} = e^{0.165\times20} = e^{3.3} \approx 27.11.}
\end{equation*}

Because gain grows exponentially with length, a 33\% increase in rod length more than doubles the small-signal gain from 12 to 27.

% ============================================================================

\begin{question}[Question 5]A 15-cm long rod of Nd:glass used as a laser amplifier has a total
small-signal gain of 10 at $\lambda_{o}=1.06\,\mu$m. Use the data in Table 13.2-1 on page 480
in Chapter 13 of Reference 1 to determine the population difference $N$ required to achieve
this gain.
\end{question}

\textbf{\large Solution --- Question 5}

From Table 13.2-1 for Nd:glass: $\sigma_o = 3\times10^{-20}$\,cm$^{2}$. The required gain coefficient for $G_o=10$ over $L=15$\,cm is:

\begin{equation*}
  10 = e^{\gamma_o\times15} \quad\Longrightarrow\quad
  \gamma_o = \frac{\ln 10}{15} = 0.154\;\text{cm}^{-1}.
\end{equation*}

Since $\gamma_o = \sigma_o N$, the required population difference is:

\begin{equation*}
  N = \frac{\gamma_o}{\sigma_o} = \frac{0.154}{3\times10^{-20}}
\end{equation*}

\begin{equation*}
  \boxed{N \approx 5.13\times10^{18}\;\text{cm}^{-3}.}
\end{equation*}

% ============================================================================

\begin{question}[Question 6]Consider a homogeneously broadened laser amplifying medium of length
$d=10$\,cm and saturation photon flux density $\phi_{s}=4\times10^{18}$\,photons/cm$^{2}$-s.
It is known that a photon flux density at the input $\Phi(0)=4\times10^{15}$\,photons/cm$^{2}$-s
produces a photon flux density at the output $\Phi(d)=4\times10^{16}$\,photons/cm$^{2}$-s.

\begin{enumerate}[label=(\alph*)]
  \item Determine the small-signal gain coefficient of the system, $G_{o}$.
  \item Determine the small-signal gain coefficient $\gamma_{o}$ of the system.
  \item What is the photon flux density at which the gain coefficient decreases by a factor of 5?
  \item Determine the gain coefficient when the input photon flux density is
        $\Phi(0)=4\times10^{19}$\,photons/cm$^{2}$-s. Under these conditions is the gain of
        the system greater than, less than, or equal to the small-signal gain determined in
        part (a)?
\end{enumerate}
\end{question}

\textbf{\large Solution --- Question 6}

\textbf{Part (a): Small-signal gain.}

Because both the input and output photon flux densities are far below the saturation flux ($\Phi(0), \Phi(d) \ll \Phi_s$), the medium operates purely in the linear (small-signal) regime. The small-signal gain $G_o$ is the ratio of output to input flux:

\begin{equation*}
  G_o = \frac{\Phi(d)}{\Phi(0)} = \frac{4\times10^{16}}{4\times10^{15}}
\end{equation*}

\begin{equation*}
  \boxed{G_o = 10.}
\end{equation*}

\textbf{Part (b): Gain coefficient.}

In the small-signal regime, the flux grows exponentially according to $G_o = e^{\gamma_o d}$. Solving for the small-signal gain coefficient $\gamma_o$:

\begin{equation*}
  \gamma_o = \frac{\ln(G_o)}{d} = \frac{\ln(10)}{10} \approx 0.230\;\text{cm}^{-1}.
\end{equation*}

\begin{equation*}
  \boxed{\gamma_o \approx 0.230\;\text{cm}^{-1}.}
\end{equation*}

\textbf{Part (c): Saturation flux condition.}

The saturated gain coefficient for a homogeneously broadened medium follows $\gamma = \gamma_o / (1 + \Phi/\Phi_s)$. We seek the flux where $\gamma = \gamma_o / 5$:

\begin{equation*}
  \frac{\gamma_o}{5} = \frac{\gamma_o}{1 + \Phi/\Phi_s} \quad\Longrightarrow\quad
  1 + \frac{\Phi}{\Phi_s} = 5 \quad\Longrightarrow\quad \Phi = 4\Phi_s.
\end{equation*}

\begin{equation*}
  \boxed{\Phi = 4(4\times10^{18}) = 1.6\times10^{19}\;\text{photons/cm}^{2}\text{-s}.}
\end{equation*}

\textbf{Part (d): Saturated gain.}

If the input flux is $\Phi(0) = 4\times10^{19}$\,photons/cm$^2$-s, the signal immediately saturates the amplifier because $\Phi(0) = 10\Phi_s$. The gain coefficient at the input drops to:

\begin{equation*}
  \gamma(0) = \frac{\gamma_o}{1 + 10} = \frac{\gamma_o}{11} \approx 0.021\;\text{cm}^{-1}.
\end{equation*}

As the signal propagates, it continues to deplete the inversion, keeping the medium deeply saturated. Because the local gain coefficient $\gamma(z)$ is everywhere much smaller than $\gamma_o$, the overall integrated gain $G$ will be dramatically \textbf{less than} the small-signal gain $G_o=10$ determined in part (a).

% ============================================================================

\begin{question}[Question 7]Consider an atomic laser amplifier system with four levels: 0 (ground
state), 1, 2, and 3. Two pumps are applied: between the ground state and level 3 at a rate
$R_{3}$, and between the ground state and level 2 at a rate $R_{2}$. Population inversion can
occur between levels 3 and 1 and/or between levels 2 and 1. Assume that no input signal is
launched to the amplifier ($\Phi=0$). Also assume that the decay from level 3 to 2 and that the
decay from levels 3 and 2 to the ground state are not possible.

\begin{enumerate}[label=(\alph*)]
  \item Write the rate equations for levels 1, 2, and 3 in terms of the lifetimes $\tau_{10}$,
        $\tau_{31}$, and $\tau_{21}$.
  \item Determine the steady state populations $N_{1}$, $N_{2}$, and $N_{3}$.
  \item Find the relation between $\tau_{10}$, $\tau_{31}$, and $\tau_{21}$ (which is
        independent of $R_{2}$ or $R_{3}$) that enables simultaneous population inversion
        between 3 and 1 and between 2 and 1.
\end{enumerate}
\end{question}

\textbf{\large Solution --- Question 7}

\textbf{Part (a): Rate equations.}

With zero input signal ($\Phi=0$), there are no stimulated transitions. Levels 3 and 2 are populated exclusively by their respective ground-state pumps, and they decay only via spontaneous emission to level 1. Level 1 subsequently decays to the ground state (level 0). The rate equations are:

\begin{align*}
  \frac{dN_3}{dt} &= R_3 - \frac{N_3}{\tau_{31}} \\[1ex]
  \frac{dN_2}{dt} &= R_2 - \frac{N_2}{\tau_{21}} \\[1ex]
  \frac{dN_1}{dt} &= \frac{N_3}{\tau_{31}} + \frac{N_2}{\tau_{21}} - \frac{N_1}{\tau_{10}}
\end{align*}

\textbf{Part (b): Steady-state populations.}

Setting the time derivatives to zero yields the steady-state populations, which perfectly trace the pumping rates and lifetimes:

\begin{align*}
  N_3 &= R_3 \tau_{31} \\[1ex]
  N_2 &= R_2 \tau_{21} \\[1ex]
  N_1 &= \tau_{10}\left( \frac{N_3}{\tau_{31}} + \frac{N_2}{\tau_{21}} \right) = \tau_{10}(R_3 + R_2)
\end{align*}

\textbf{Part (c): Condition for simultaneous inversion.}

Simultaneous inversion requires both $N_3 > N_1$ and $N_2 > N_1$. Substituting our steady-state expressions:

\begin{align*}
  R_3 \tau_{31} &> \tau_{10}(R_3 + R_2) \quad\Longrightarrow\quad \frac{\tau_{31}}{\tau_{10}} > 1 + \frac{R_2}{R_3} \\[1ex]
  R_2 \tau_{21} &> \tau_{10}(R_3 + R_2) \quad\Longrightarrow\quad \frac{\tau_{21}}{\tau_{10}} > 1 + \frac{R_3}{R_2}
\end{align*}

To find a general condition independent of the pump intensities $R_2$ and $R_3$, we invert both inequalities:

\begin{equation*}
  \frac{\tau_{10}}{\tau_{31}} < \frac{R_3}{R_3 + R_2} \qquad\text{and}\qquad \frac{\tau_{10}}{\tau_{21}} < \frac{R_2}{R_3 + R_2}.
\end{equation*}

Summing these inequalities eliminates the pump ratios completely:

\begin{equation*}
  \frac{\tau_{10}}{\tau_{31}} + \frac{\tau_{10}}{\tau_{21}} < 1 \quad\Longrightarrow\quad
  \boxed{\frac{1}{\tau_{31}} + \frac{1}{\tau_{21}} < \frac{1}{\tau_{10}}.}
\end{equation*}

This elegantly demonstrates that for simultaneous dual-line lasing, the combined spontaneous decay rates from the upper levels must not exceed the clearance rate from the shared lower level; otherwise, atoms accumulate in level 1 and destroy the inversion.

% ============================================================================

\begin{question}[Question 8]Consider the ideal four-level laser system shown in the figure. The pump
excites atoms from level 0 to level 3 at a rate $R$, then the atoms decay almost instantly from
level 3 to level 2. The spontaneous decay rates as well as the absorption and stimulated emission
rates are illustrated in the figure. The energies are $E_{1}=0.2$\,eV, $E_{2}=1.365$\,eV,
$E_{3}=2.1$\,eV with respect to the ground level. The lifetimes are $\tau_{1}=30$\,ns and
$\tau_{2}=1.2$\,ms, the linewidth is $\Delta\nu=120$\,GHz, and the refractive index of the
laser medium is $n=1.82$.

\begin{enumerate}[label=(\alph*)]
  \item Write down the rate equations for the atoms at levels 2 and 1.
  \item Show that the steady state solution of these equations gives a population density
        difference $N_{2}-N_{1}=R(\tau_{2}-\tau_{1})/(1+\tau_{2}W_{i})$.
  \item Calculate the transition cross section (in cm$^{2}$) at
        $\nu_{o}(=(E_{2}-E_{1})/h)$.
  \item Show that the gain coefficient may be written in the form
        $\gamma(\nu)=R(\tau_{2}-\tau_{1})\sigma(\nu)/(1+\tau_{2}\sigma(\nu)\Phi)$.
  \item Calculate the saturation photon-flux density $\phi_{s}$ (in photons/cm$^{2}$-s) at
        $\nu_{o}$.
  \item Calculate the photon flux density at which the gain decreases by a factor of 5.
  \item If the length of the laser medium is 10\,cm, find the pump power (in W/cm$^{3}$)
        required to achieve a small-signal amplifier gain of 15\,dB (at $\nu=\nu_{o}$).
\end{enumerate}
\end{question}

\textbf{\large Solution --- Question 8}

\textbf{Part (a): Rate equations.}

Because atoms decay almost instantly from level 3 to 2, the pump $R$ effectively populates level 2 directly. Incorporating the stimulated emission rate $W_i = \sigma(\nu)\Phi$:

\begin{align*}
  \frac{dN_2}{dt} &= R - W_i(N_2 - N_1) - \frac{N_2}{\tau_2} \\[1ex]
  \frac{dN_1}{dt} &= W_i(N_2 - N_1) + \frac{N_2}{\tau_2} - \frac{N_1}{\tau_1}
\end{align*}

\textbf{Part (b): Steady-state inversion.}

Summing the two rate equations at steady state ($d/dt = 0$) cancels out the inter-level transitions, revealing the strict dependency of the lower level on the pump:

\begin{equation*}
  R - \frac{N_1}{\tau_1} = 0 \quad\Longrightarrow\quad N_1 = R\tau_1.
\end{equation*}

We seek the population difference $\Delta N = N_2 - N_1$. Rewriting $N_2 = \Delta N + R\tau_1$ and substituting into the first rate equation:

\begin{equation*}
  0 = R - W_i\Delta N - \frac{\Delta N + R\tau_1}{\tau_2}.
\end{equation*}

Multiplying by $\tau_2$ and isolating $\Delta N$:

\begin{equation*}
  R\tau_2 - W_i\tau_2\Delta N - \Delta N - R\tau_1 = 0 \quad\Longrightarrow\quad
  \Delta N(1 + \tau_2 W_i) = R(\tau_2 - \tau_1).
\end{equation*}

\begin{equation*}
  \boxed{N_2 - N_1 = \frac{R(\tau_2 - \tau_1)}{1 + \tau_2 W_i}.}
\end{equation*}

\textbf{Part (c): Transition cross section.}

The transition frequency is $\nu_o = (E_2 - E_1)/h$. Converting the energy difference:

\begin{equation*}
  \Delta E = 1.365 - 0.2 = 1.165\,\text{eV} = 1.866\times10^{-19}\;\text{J} \quad\Longrightarrow\quad
  \nu_o = 2.817\times10^{14}\;\text{Hz}.
\end{equation*}

Assuming a homogeneously broadened lineshape peaked at $g(\nu_o) \approx 1/\Delta\nu$ (using the effective linewidth approximation) and taking $\tau_\text{sp} \approx \tau_2$:

\begin{equation*}
  \sigma(\nu_o) = \frac{c^{2}}{8\pi n^{2}\nu_o^{2}\tau_2}\cdot\frac{1}{\Delta\nu}
  = \frac{(3\times10^{8})^{2}}{8\pi(1.82)^{2}(2.817\times10^{14})^{2}(1.2\times10^{-3})(120\times10^{9})}.
\end{equation*}

\begin{equation*}
  \boxed{\sigma(\nu_o) \approx 9.47\times10^{-20}\;\text{cm}^{2}.}
\end{equation*}
\emph{Note: If the strict Lorentzian peak $g(\nu_o)=2/(\pi\Delta\nu)$ is used, the cross section scales by $2/\pi$, yielding $\approx 6.03\times10^{-20}$\,cm$^{2}$.}

\textbf{Part (d): Gain coefficient formula.}

The gain coefficient is $\gamma(\nu) = \sigma(\nu)\Delta N$. Substituting $\Delta N$ from part (b) and $W_i = \sigma(\nu)\Phi$:

\begin{equation*}
  \boxed{\gamma(\nu) = \frac{R(\tau_2 - \tau_1)\sigma(\nu)}{1 + \tau_2\sigma(\nu)\Phi}.}
\end{equation*}

\textbf{Part (e): Saturation flux.}

By inspecting the denominator of the gain coefficient, the saturation photon-flux density is defined when $\tau_2\sigma(\nu_o)\Phi = 1$:

\begin{equation*}
  \Phi_s = \frac{1}{\tau_2\sigma(\nu_o)} = \frac{1}{(1.2\times10^{-3})(9.47\times10^{-20})}.
\end{equation*}

\begin{equation*}
  \boxed{\Phi_s \approx 8.80\times10^{21}\;\text{photons/cm}^{2}\text{-s}.}
\end{equation*}

\textbf{Part (f): Gain decrease condition.}

For the gain to decrease by a factor of 5, the denominator must equal 5:

\begin{equation*}
  1 + \frac{\Phi}{\Phi_s} = 5 \quad\Longrightarrow\quad \Phi = 4\Phi_s.
\end{equation*}

\begin{equation*}
  \boxed{\Phi = 4(8.80\times10^{21}) = 3.52\times10^{22}\;\text{photons/cm}^{2}\text{-s}.}
\end{equation*}

\textbf{Part (g): Required pump power.}

A small-signal gain of 15\,dB over $L=10$\,cm translates to a linear gain $G_o$:

\begin{equation*}
  10\log_{10}(G_o) = 15 \quad\Longrightarrow\quad G_o = 10^{1.5} \approx 31.62.
\end{equation*}

The small-signal gain coefficient is $\gamma_o = \ln(31.62)/10 \approx 0.345$\,cm$^{-1}$.
Using the unsaturated formula $\gamma_o = R(\tau_2 - \tau_1)\sigma(\nu_o)$, we find the required pump rate:

\begin{equation*}
  R = \frac{\gamma_o}{(\tau_2 - \tau_1)\sigma(\nu_o)}
  = \frac{0.345}{(1.2\times10^{-3} - 30\times10^{-9})(9.47\times10^{-20})}
  \approx 3.04\times10^{21}\;\text{cm}^{-3}\text{s}^{-1}.
\end{equation*}

Since each pumping event requires an energy $E_3 = 2.1$\,eV $= 3.36\times10^{-19}$\,J, the volumetric pump power is:

\begin{equation*}
  P = R E_3 = (3.04\times10^{21})(3.36\times10^{-19}).
\end{equation*}

\begin{equation*}
  \boxed{P \approx 1021\;\text{W/cm}^{3}.}
\end{equation*}

% ============================================================================

\begin{question}[Formative Question]Consider the energy-level system of a proposed laser
as shown in the figure.
\begin{enumerate}[label=(\alph*)]
  \item Find the equilibrium populations of states 1 and 2 when $R=10^{20}$\,cm$^{-3}$/sec.
        Assume that $\tau_{2}=20$\,ns and $\tau_{1}=1\,\mu$s and neglect the depletion of
        state 0.
  \item Suppose that $R(t)=R_{o}H(t)$, where $H(t)$ is the Heaviside step function. Plot the
        time variation of the gain coefficient assuming that $\sigma=10^{-12}$\,cm$^{2}$ and
        that $N_{2}(0)=0$ (neglect stimulated emission).
\end{enumerate}
\end{question}

\textbf{\large Solution --- Formative Question}

\textbf{Part (a): Equilibrium populations.}

At steady state, the rate equations yield $N_2 = R\tau_2$ and $N_1 = R\tau_1$:

\begin{align*}
  N_2 &= (10^{20})(20\times10^{-9}) = 2\times10^{12}\;\text{cm}^{-3} \\[1ex]
  N_1 &= (10^{20})(1\times10^{-6}) = 1\times10^{14}\;\text{cm}^{-3}
\end{align*}

Because $\tau_1 \gg \tau_2$, atoms bottleneck in the lower level. The inversion $\Delta N = N_2 - N_1 = -9.8\times10^{13}$\,cm$^{-3}$ is highly negative, meaning the medium acts as an absorber, not an amplifier.

\textbf{Part (b): Time-varying gain coefficient.}

With a step-function pump $R(t) = R_o H(t)$ turning on at $t=0$, we solve the transient rate equations:

\begin{align*}
  \frac{dN_2}{dt} &= R_o - \frac{N_2}{\tau_2} \quad\Longrightarrow\quad N_2(t) = R_o\tau_2(1 - e^{-t/\tau_2}) \\[1ex]
  \frac{dN_1}{dt} &= \frac{N_2}{\tau_2} - \frac{N_1}{\tau_1} = R_o(1 - e^{-t/\tau_2}) - \frac{N_1}{\tau_1}
\end{align*}

Solving the first-order linear differential equation for $N_1(t)$ with $N_1(0)=0$:

\begin{equation*}
  N_1(t) = R_o\tau_1 \left( 1 - \frac{\tau_1 e^{-t/\tau_1} - \tau_2 e^{-t/\tau_2}}{\tau_1 - \tau_2} \right)
\end{equation*}

The transient gain coefficient is $\gamma(t) = \sigma [N_2(t) - N_1(t)]$. At early times ($t \sim \tau_2$), $N_2$ fills rapidly while $N_1$ is still empty, creating a brief window of positive transient inversion (and thus positive gain). However, as $t \to \tau_1$, $N_1$ overtakes $N_2$, and the gain coefficient plummets below zero, settling at the highly absorptive steady-state value derived in part (a).
This system inherently fails as a CW laser but could hypothetically function as a pulsed laser if operated on timescales shorter than $\tau_1$.
